% ============================================================
% CHAPITRE 3 : MÉTHODOLOGIE
% ============================================================

\chapter{Méthodologie de Développement}

\section{Introduction}

Ce chapitre présente la méthodologie adoptée pour le développement et la validation du système DMS. Le projet suit rigoureusement le \keyword{cycle en V} (\textit{V-Model}), méthodologie standard dans l'industrie automobile, complétée par les approches de test et validation \keyword{SIL}, \keyword{MIL} et \keyword{HIL}.

\section{Le Cycle en V Automobile}

\subsection{Présentation du Modèle}

Le \keyword{cycle en V} est la méthodologie de développement de référence dans l'industrie automobile, préconisée par la norme ISO 26262 \cite{iso26262}. Il organise le développement en deux branches symétriques :

\begin{itemize}[leftmargin=2cm]
    \item \textbf{Branche descendante (gauche)} : Phases de spécification et de conception, du niveau système au niveau composant.
    \item \textbf{Branche ascendante (droite)} : Phases de test et de validation, du niveau composant au niveau système.
\end{itemize}

Chaque phase de conception sur la branche gauche est directement associée à une phase de test correspondante sur la branche droite.

\begin{figure}[H]
\centering
\begin{tikzpicture}[
    box/.style={rectangle, draw, fill=#1, text width=3.5cm, text centered, minimum height=1cm, rounded corners=2pt, font=\small\bfseries},
    arrow/.style={->, >=stealth, thick},
    dashed_arrow/.style={->, >=stealth, thick, dashed, color=red!60!black}
]
    % Branche gauche (descendante)
    \node[box=stellantisblue!15] (req) at (0, 6) {Spécifications\\Système};
    \node[box=stellantisblue!20] (arch) at (-2, 4.5) {Architecture\\Système};
    \node[box=stellantisblue!25] (design) at (-4, 3) {Conception\\Détaillée};
    \node[box=stellantisblue!30] (impl) at (-5, 1.5) {Implémentation\\(Codage)};
    
    % Point bas
    \node[box=expleored!20] (unit) at (-3, 0) {Tests\\Unitaires};
    
    % Branche droite (ascendante)
    \node[box=green!15] (integ) at (0, 1.5) {Tests\\d'Intégration};
    \node[box=green!20] (sysv) at (2, 3) {Validation\\Système};
    \node[box=green!25] (accept) at (4, 4.5) {Tests\\d'Acceptation};
    \node[box=green!30] (valid) at (2, 6) {Validation\\Finale};
    
    % Flèches descendantes
    \draw[arrow] (req) -- (arch);
    \draw[arrow] (arch) -- (design);
    \draw[arrow] (design) -- (impl);
    \draw[arrow] (impl) -- (unit);
    
    % Flèches ascendantes
    \draw[arrow] (unit) -- (integ);
    \draw[arrow] (integ) -- (sysv);
    \draw[arrow] (sysv) -- (accept);
    \draw[arrow] (accept) -- (valid);
    
    % Correspondances horizontales
    \draw[dashed_arrow] (req) -- node[above, font=\tiny, color=red!60!black] {vérifie} (valid);
    \draw[dashed_arrow] (arch) -- node[above, font=\tiny, color=red!60!black] {vérifie} (accept);
    \draw[dashed_arrow] (design) -- node[above, font=\tiny, color=red!60!black] {vérifie} (sysv);
    \draw[dashed_arrow] (impl) -- node[above, font=\tiny, color=red!60!black] {vérifie} (integ);
    
    % Labels des niveaux
    \node[font=\tiny\itshape, color=gray] at (-6.5, 4.5) {Niveau Système};
    \node[font=\tiny\itshape, color=gray] at (-6.5, 3) {Niveau Module};
    \node[font=\tiny\itshape, color=gray] at (-6.5, 1.5) {Niveau Composant};
    
\end{tikzpicture}
\caption{Le cycle en V appliqué au projet DMS.}
\label{fig:cycle_v}
\end{figure}

\subsection{Application au Projet DMS}

Le tableau~\ref{tab:cycle_v_dms} détaille l'application du cycle en V à notre projet :

\begin{table}[H]
\centering
\caption{Application du cycle en V au projet DMS.}
\label{tab:cycle_v_dms}
\renewcommand{\arraystretch}{1.4}
\begin{tabularx}{\textwidth}{|>{\bfseries\small}p{2.5cm}|X|X|}
\hline
\rowcolor{stellantisblue!15}
\textbf{Phase} & \textbf{Activités (Branche gauche)} & \textbf{Activités (Branche droite)} \\
\hline
Spécifications & Définition des exigences fonctionnelles et non-fonctionnelles du DMS (PERCLOS, alertes, temps de réponse) & Validation finale : le système répond aux exigences initiales \\
\hline
Architecture & Décomposition en modules (acquisition, détection, analyse, ECU), définition des interfaces & Tests d'acceptation : validation de chaque module selon les spécifications \\
\hline
Conception détaillée & Architecture CNN, algorithmes de détection, protocoles de communication & Validation système : tests d'intégration multi-modules \\
\hline
Implémentation & Codage Python/PyTorch, modèles Simulink, firmware robot & Tests unitaires : validation de chaque fonction/classe \\
\hline
\end{tabularx}
\end{table}

\section{Approches de Test et Validation}

La validation du système DMS suit une progression en trois niveaux d'abstraction, conformément aux pratiques de l'industrie automobile \cite{expleo2024} :

\subsection{MIL -- Model-in-the-Loop}

Le test \keyword{MIL} (\textit{Model-in-the-Loop}) est la première phase de validation. Elle consiste à tester les algorithmes et les modèles mathématiques dans un environnement de simulation pure.

\begin{figure}[H]
\centering
\begin{tikzpicture}[
    block/.style={rectangle, draw, fill=#1, text width=3cm, text centered, minimum height=1.2cm, rounded corners=3pt, font=\small},
    arrow/.style={->, >=stealth, thick}
]
    \node[block=orange!15] (model) {Modèle\\Algorithmique\\(MATLAB/Simulink)};
    \node[block=blue!15, right=2cm of model] (plant) {Modèle de\\l'Environnement\\(Plant Model)};
    \node[block=green!15, below=1.5cm of $(model)!0.5!(plant)$] (result) {Résultats\\de Simulation};
    
    \draw[arrow] (model) -- node[above, font=\tiny] {signaux} (plant);
    \draw[arrow] (plant) -- node[above, font=\tiny] {feedback} (model);
    \draw[arrow] (model) -- (result);
    \draw[arrow] (plant) -- (result);
    
    \node[draw, dashed, fit=(model)(plant), inner sep=0.5cm, label=above:{\textbf{Simulation Pure (PC)}}] {};
\end{tikzpicture}
\caption{Architecture de test MIL.}
\label{fig:mil}
\end{figure}

\textbf{Dans notre projet :}
\begin{itemize}
    \item Simulation du système DMS complet sous MATLAB/Simulink
    \item Génération de signaux synthétiques (état des yeux, angles de tête)
    \item Validation des algorithmes de calcul de PERCLOS et de scores
    \item Vérification de la logique ECU avec la machine à états Stateflow
\end{itemize}

\subsection{SIL -- Software-in-the-Loop}

Le test \keyword{SIL} (\textit{Software-in-the-Loop}) intègre le code logiciel réel dans la boucle de simulation, tout en utilisant un environnement simulé pour les entrées/sorties.

\begin{figure}[H]
\centering
\begin{tikzpicture}[
    block/.style={rectangle, draw, fill=#1, text width=3cm, text centered, minimum height=1.2cm, rounded corners=3pt, font=\small},
    arrow/.style={->, >=stealth, thick}
]
    \node[block=red!15] (code) {Code Logiciel\\Réel\\(Python/PyTorch)};
    \node[block=blue!15, right=2cm of code] (sim) {Environnement\\Simulé\\(Images de test)};
    \node[block=green!15, below=1.5cm of $(code)!0.5!(sim)$] (result) {Résultats\\de Test};
    
    \draw[arrow] (sim) -- node[above, font=\tiny] {images} (code);
    \draw[arrow] (code) -- node[below, font=\tiny] {décisions} (sim);
    \draw[arrow] (code) -- (result);
    
    \node[draw, dashed, fit=(code)(sim), inner sep=0.5cm, label=above:{\textbf{Logiciel Réel + Simulation (PC)}}] {};
\end{tikzpicture}
\caption{Architecture de test SIL.}
\label{fig:sil}
\end{figure}

\textbf{Dans notre projet :}
\begin{itemize}
    \item Exécution du code Python réel avec des images de test synthétiques et réelles
    \item Validation des modèles CNN (EyeStateCNN, FaceDetectionCNN)
    \item Tests des scénarios complets (conducteur attentif, fatigué, distrait, mixte)
    \item 34 tests unitaires couvrant l'ensemble des modules
\end{itemize}

\subsection{HIL -- Hardware-in-the-Loop}

Le test \keyword{HIL} (\textit{Hardware-in-the-Loop}) intègre le matériel réel (ECU, capteurs, actionneurs) dans la boucle de test, simulant uniquement l'environnement extérieur.

\begin{figure}[H]
\centering
\begin{tikzpicture}[
    block/.style={rectangle, draw, fill=#1, text width=2.8cm, text centered, minimum height=1.2cm, rounded corners=3pt, font=\small},
    arrow/.style={->, >=stealth, thick}
]
    \node[block=red!20] (ecu) {ECU Réel\\(Raspberry Pi)};
    \node[block=orange!15, above right=0.8cm and 2cm of ecu] (cam) {Caméra\\Réelle};
    \node[block=blue!15, right=2cm of ecu] (sim) {Environnement\\Simulé};
    \node[block=green!15, below right=0.8cm and 2cm of ecu] (act) {Actionneurs\\Réels};
    
    \draw[arrow] (cam) -- node[above, font=\tiny, sloped] {images} (ecu);
    \draw[arrow] (sim) -- node[above, font=\tiny] {scénarios} (ecu);
    \draw[arrow] (ecu) -- node[below, font=\tiny, sloped] {commandes} (act);
    
    \node[draw, dashed, fit=(ecu)(cam)(act), inner sep=0.5cm, label=below:{\textbf{Matériel Réel + Simulation Environnement}}] {};
\end{tikzpicture}
\caption{Architecture de test HIL.}
\label{fig:hil}
\end{figure}

\textbf{Dans notre projet :}
\begin{itemize}
    \item Prototype robot avec Raspberry Pi comme ECU
    \item Caméra USB réelle pour l'acquisition vidéo
    \item Actionneurs (LED, buzzer, servo-moteurs) simulant les alertes
    \item Communication CAN/UART entre les composants
    \item Tests des scénarios en conditions réelles
\end{itemize}

\subsection{Comparaison des Approches}

\begin{table}[H]
\centering
\caption{Comparaison des approches de test MIL, SIL et HIL.}
\label{tab:mil_sil_hil}
\renewcommand{\arraystretch}{1.4}
\begin{tabularx}{\textwidth}{|>{\bfseries}l|X|X|X|}
\hline
\rowcolor{stellantisblue!15}
\textbf{Critère} & \textbf{MIL} & \textbf{SIL} & \textbf{HIL} \\
\hline
Logiciel & Modèle (Simulink) & Code réel (Python) & Code réel embarqué \\
\hline
Matériel & Aucun (simulation) & Aucun (PC) & ECU + capteurs réels \\
\hline
Environnement & Simulé & Simulé & Partiellement simulé \\
\hline
Coût & Faible & Moyen & Élevé \\
\hline
Fidélité & Basse & Moyenne & Haute \\
\hline
Temps de test & Rapide & Moyen & Lent \\
\hline
Phase du cycle en V & Conception & Implémentation & Intégration \\
\hline
\end{tabularx}
\end{table}

\section{Outils de Test et Validation}

\subsection{Outils Logiciels}

Les outils suivants sont utilisés pour le test et la validation :

\begin{table}[H]
\centering
\caption{Outils de test et validation utilisés.}
\label{tab:outils_test}
\renewcommand{\arraystretch}{1.3}
\begin{tabularx}{\textwidth}{|l|l|X|}
\hline
\rowcolor{stellantisblue!15}
\textbf{Outil} & \textbf{Type} & \textbf{Utilisation} \\
\hline
Python unittest & Test unitaire & 34 tests unitaires pour tous les modules \\
\hline
MATLAB/Simulink & Simulation & Modélisation et simulation MIL du système DMS \\
\hline
Stateflow & Machine à états & Modélisation de la logique ECU \\
\hline
Simulink Test & Test automatisé & Tests de régression des modèles Simulink \\
\hline
OpenCV & Vision par ordinateur & Traitement d'images et détection \\
\hline
PyTorch & Deep Learning & Entraînement et inférence des CNN \\
\hline
\end{tabularx}
\end{table}

\subsection{Métriques de Validation}

Les métriques suivantes sont utilisées pour évaluer les performances du système :

\begin{enumerate}[leftmargin=2cm]
    \item \textbf{Précision de détection} : $\text{Précision} = \frac{VP}{VP + FP}$
    \item \textbf{Rappel} : $\text{Rappel} = \frac{VP}{VP + FN}$
    \item \textbf{F1-Score} : $F_1 = 2 \cdot \frac{\text{Précision} \cdot \text{Rappel}}{\text{Précision} + \text{Rappel}}$
    \item \textbf{Temps de réponse} : Latence entre la détection et la génération d'alerte
    \item \textbf{Taux de faux positifs} : Alertes générées sans danger réel
    \item \textbf{Taux de détection} : Pourcentage de situations dangereuses correctement identifiées
\end{enumerate}

\section{Gestion de Projet}

\subsection{Planning Prévisionnel}

Le projet est planifié sur une durée de 6 mois, décomposé en sprints de 2 semaines :

\begin{table}[H]
\centering
\caption{Planning prévisionnel du projet DMS.}
\label{tab:planning}
\renewcommand{\arraystretch}{1.3}
\begin{tabularx}{\textwidth}{|c|l|X|}
\hline
\rowcolor{stellantisblue!15}
\textbf{Mois} & \textbf{Phase} & \textbf{Livrables} \\
\hline
1 & Analyse et spécification & Cahier des charges, analyse QQOQCP, spécifications fonctionnelles \\
\hline
2 & Conception & Architecture système, conception CNN, modèles Simulink \\
\hline
3 & Implémentation (Partie 1) & Code Python, modèles CNN entraînés, algorithmes de détection \\
\hline
4 & Implémentation (Partie 2) & Logique ECU, analyse comportementale, prototype robot \\
\hline
5 & Test et Validation & Tests unitaires, SIL, MIL, HIL, scénarios de validation \\
\hline
6 & Rédaction et soutenance & Rapport PFE, préparation soutenance, démonstration \\
\hline
\end{tabularx}
\end{table}

\section{Synthèse}

La méthodologie adoptée pour ce projet combine :
\begin{itemize}
    \item Le \textbf{cycle en V} pour structurer le développement de manière rigoureuse
    \item Les approches \textbf{MIL}, \textbf{SIL} et \textbf{HIL} pour une validation progressive et exhaustive
    \item Des outils industriels standard (MATLAB/Simulink, Python, PyTorch)
    \item Des métriques de performance quantitatives et objectives
\end{itemize}

Cette approche garantit la qualité, la fiabilité et la traçabilité du système développé, conformément aux exigences de l'industrie automobile et aux standards de Stellantis.
